\documentclass[11pt]{amsart}

\usepackage[T1]{fontenc}
\usepackage{lmodern}
\usepackage{microtype}
\usepackage{amsmath,amssymb,amsthm}
\usepackage[hidelinks]{hyperref}

\newtheorem{theorem}{Theorem}
\newtheorem{lemma}[theorem]{Lemma}
\newtheorem{corollary}[theorem]{Corollary}
\newtheorem{proposition}[theorem]{Proposition}
\theoremstyle{remark}
\newtheorem{remark}[theorem]{Remark}

\newcommand{\R}{\mathbb{R}}
\newcommand{\Z}{\mathbb{Z}}
\newcommand{\Torus}{\mathbb{T}}
\newcommand{\OO}{\operatorname{O}}
\newcommand{\SO}{\operatorname{SO}}
\newcommand{\Imag}{\operatorname{Im}}

\title[Maximum Binary Case of Ribot's Zariski-Closedness]{The Maximum Binary Case of Ribot's Zariski-Closedness Conjecture}
\date{}

\subjclass[2020]{Primary 14P05; Secondary 15A69, 14L30, 22C05, 05B15}
\keywords{real Zariski closure, orbit of the orthogonal group, tensor,
partial Latin hyperrectangle, parity class, character duality, compact torus}

\hypersetup{
  pdftitle={The Maximum Binary Case of Ribot's Zariski-Closedness Conjecture}
}

\begin{document}

\begin{abstract}
Ribot conjectured that the orthogonal orbit of the coordinate subspace
supported on any partial Latin hyperrectangle is real Zariski closed. We
prove the conjecture for every maximum-cardinality support in binary format.
Such supports are exactly the two parity classes in $\{0,1\}^d$. A Fourier
change of coordinates diagonalizes the action of $\SO(2)^d$, and character
duality for compact tori gives explicit polynomial equations indexed by the
integer relations among the resulting characters. In format
$2\times2\times2$, the relation lattice is cyclic and the resulting quartic
is a nonzero scalar multiple of Ribot's quartic for the seven-dimensional
component. Thus the orbit itself, not merely its Zariski closure, is that
hypersurface.
\end{abstract}

\maketitle

\section{Statement}

For $x=(x_1,\ldots,x_d)\in\{0,1\}^d$, write
\[
e_x=e_{x_1}\otimes\cdots\otimes e_{x_d}.
\]
Given $L\subseteq\{0,1\}^d$, set
\[
V_L=\operatorname{span}_{\R}\{e_x:x\in L\},
\qquad
X_L=\OO(2)^d\cdot V_L\subseteq(\R^2)^{\otimes d}.
\]
A partial Latin hyperrectangle is a set whose distinct elements have Hamming
distance at least two. Ribot conjectured that $X_L$ is Zariski closed for
every such support~\cite[Conjecture~4.4]{Ribot}. Known positive cases include
matrices and orthogonally decomposable tensors~\cite{BDHR,Ribot}.

For $\delta\in\{0,1\}$, define the parity support
\[
E_\delta=\{x\in\{0,1\}^d:|x|\equiv\delta\pmod 2\}.
\]
Choose $\Omega\subseteq\{\pm1\}^d$ containing one vector from each antipodal
pair. For a real tensor $T=\sum_x t_xe_x$, put
\[
p_\varepsilon(T)
=2^{-d}\sum_{x\in\{0,1\}^d}t_x i^{|x|}\varepsilon^x,
\qquad
\varepsilon^x=\prod_{k=1}^d\varepsilon_k^{x_k}.
\]
Let
\[
A:\Z^\Omega\longrightarrow\Z^d,
\qquad
A(m)=\sum_{\varepsilon\in\Omega}m_\varepsilon\varepsilon,
\qquad
\Lambda=\ker A,
\]
and write $m_\varepsilon^+=\max\{m_\varepsilon,0\}$ and
$m_\varepsilon^-=\max\{-m_\varepsilon,0\}$. For $m\in\Lambda$, define
\[
\Phi_m(T)=
\Imag\!\left(
\prod_{\varepsilon\in\Omega}
p_\varepsilon(T)^{m_\varepsilon^+}
\overline{p_\varepsilon(T)}^{m_\varepsilon^-}
\right).
\]
Each $\Phi_m$ is a real homogeneous polynomial in the entries of $T$.

\begin{theorem}\label{thm:main}
For every $d\ge2$,
\[
X_{E_0}=X_{E_1}
=\{T:\Phi_m(T)=0\text{ for every }m\in\Lambda\}.
\]
Consequently, $X_{E_0}$ and $X_{E_1}$ are real Zariski closed. Moreover,
\[
\operatorname{rank}\Lambda=2^{d-1}-d.
\]
\end{theorem}

\begin{corollary}\label{cor:maximum}
Ribot's Conjecture~4.4 holds for every maximum-cardinality partial Latin
hyperrectangle in $\{0,1\}^d$.
\end{corollary}

\begin{remark}
Combining Corollary~\ref{cor:maximum} with Ribot's dimension formula and
component classification~\cite[Lemma~4.5 and Theorem~1.7]{Ribot} shows that
$X_{E_0}$ is the unique top-dimensional irreducible component of Ribot's
bato variety in $(\R^2)^{\otimes d}$; its dimension is $2^{d-1}+d$.
\end{remark}

\begin{remark}
Although Theorem~\ref{thm:main} is stated using all $m\in\Lambda$, finitely
many of the equations $\Phi_m=0$ suffice, since the real polynomial ring in
the entries of $T$ is Noetherian.
\end{remark}

\section{Proof}

Every element of $\OO(2)$ has the form $R_\theta D^a$, where
$R_\theta\in\SO(2)$, $a\in\{0,1\}$, and $D=\operatorname{diag}(1,-1)$.
Since $D$ preserves every coordinate subspace $V_L$,
\begin{equation}\label{eq:SO-reduction}
X_L=\SO(2)^d\cdot V_L.
\end{equation}
Swapping $e_0$ and $e_1$ in one factor sends $E_0$ to $E_1$, so
$X_{E_0}=X_{E_1}$.

\begin{lemma}\label{lem:fourier}
For $T\in(\R^2)^{\otimes d}$ and $\varepsilon\in\{\pm1\}^d$,
\[
p_{-\varepsilon}(T)=\overline{p_\varepsilon(T)}.
\]
If $\theta=(\theta_1,\ldots,\theta_d)$, then
\[
p_\varepsilon\bigl((R_{\theta_1},\ldots,R_{\theta_d})\cdot T\bigr)
=e^{i\langle\varepsilon,\theta\rangle}p_\varepsilon(T).
\]
Finally,
\[
T\in V_{E_0}
\quad\Longleftrightarrow\quad
p_\varepsilon(T)\in\R
\text{ for every }\varepsilon\in\Omega.
\]
\end{lemma}

\begin{proof}
Set $f_\sigma=e_0-i\sigma e_1$ for $\sigma\in\{\pm1\}$. Then
$R_\theta f_\sigma=e^{i\sigma\theta}f_\sigma$, and the coefficient of
$f_{\varepsilon_1}\otimes\cdots\otimes f_{\varepsilon_d}$ in $T$ is
$p_\varepsilon(T)$. This proves the first two assertions.

If $T\in V_{E_0}$, the defining sum for $p_\varepsilon(T)$ is real.
Conversely, if the representative coordinates are real, then
$p_{-\varepsilon}=p_\varepsilon$ for every $\varepsilon$. Fourier inversion
gives
\[
t_x=(-i)^{|x|}
\sum_{\varepsilon\in\{\pm1\}^d}p_\varepsilon(T)\varepsilon^x.
\]
Pairing $\varepsilon$ with $-\varepsilon$ makes the sum zero when $|x|$ is
odd. Hence $T\in V_{E_0}$.
\end{proof}

\begin{lemma}[Character duality]\label{lem:torus}
Let $a_1,\ldots,a_N\in\Z^d$, define
\[
B:\Z^N\longrightarrow\Z^d,
\qquad B(m)=\sum_{j=1}^Nm_ja_j,
\]
and let
\[
\phi:\Torus^d\longrightarrow\Torus^N,
\qquad
\phi(z)_j=z^{a_j}=\prod_{k=1}^dz_k^{(a_j)_k},
\]
where $\Torus=S^1$. Then
\[
\operatorname{im}\phi
=
\{u\in\Torus^N:u^m=1\text{ for every }m\in\ker B\},
\qquad
u^m=\prod_{j=1}^Nu_j^{m_j}.
\]
\end{lemma}

\begin{proof}
The forward inclusion is immediate. For the reverse inclusion, put the
integer matrix of $B$ into Smith normal form. The associated unimodular
changes of coordinates induce automorphisms of the source and target tori,
reducing $\phi$ to
\[
(z_1,\ldots,z_d)
\longmapsto
(z_1^{s_1},\ldots,z_r^{s_r},1,\ldots,1).
\]
Every power map on $S^1$ is surjective, so this image is exactly the
annihilator of the kernel of the diagonal integer map.
\end{proof}

\begin{proof}[Proof of Theorem~\ref{thm:main}]
Identify $z=(e^{i\theta_1},\ldots,e^{i\theta_d})\in\Torus^d$ with
$(R_{\theta_1},\ldots,R_{\theta_d})\in\SO(2)^d$. By
\eqref{eq:SO-reduction} and Lemma~\ref{lem:fourier}, a tensor $T$ belongs to
$X_{E_0}$ if and only if there are $z\in\Torus^d$ and
$a_\varepsilon\in\R$ such that
\begin{equation}\label{eq:phase-form}
p_\varepsilon(T)=z^\varepsilon a_\varepsilon
\qquad(\varepsilon\in\Omega).
\end{equation}
If \eqref{eq:phase-form} holds and $m\in\Lambda$, then the torus factor in
the monomial defining $\Phi_m$ is $z^{A(m)}=1$; hence $\Phi_m(T)=0$.

Conversely, suppose that $\Phi_m(T)=0$ for every $m\in\Lambda$. Let
\[
S=\{\varepsilon\in\Omega:p_\varepsilon(T)\ne0\},
\qquad
u_\varepsilon=\frac{p_\varepsilon(T)}{|p_\varepsilon(T)|}
\quad(\varepsilon\in S).
\]
For every $m\in\ker(A|_{\Z^S})$, extend $m$ by zero outside $S$. The
corresponding nonzero monomial has phase $u^m$, so its imaginary part
vanishes only if $u^m\in\{\pm1\}$. Thus
$q=(u_\varepsilon^2)_{\varepsilon\in S}$ satisfies $q^m=1$ for every such
$m$. Lemma~\ref{lem:torus} gives $w\in\Torus^d$ with
$q_\varepsilon=w^\varepsilon$. Choose $z\in\Torus^d$ coordinatewise so that
$z^2=w$. Then
\[
\left(\frac{u_\varepsilon}{z^\varepsilon}\right)^2=1,
\]
so $p_\varepsilon(T)/z^\varepsilon$ is real for $\varepsilon\in S$ and is
zero for $\varepsilon\notin S$. Lemma~\ref{lem:fourier} yields
$z^{-1}\cdot T\in V_{E_0}$, hence $T\in X_{E_0}$.

The common zero set of the real polynomials $\Phi_m$ is Zariski closed.
Finally, $|\Omega|=2^{d-1}$. The sign vectors span $\mathbb{Q}^d$: if
$v=(1,\ldots,1)$ and $v_j$ is obtained from $v$ by flipping its $j$th
coordinate, then the $j$th standard basis vector is $(v-v_j)/2$. Choosing
one vector from each antipodal pair does not change the span. Therefore
$A$ has rank $d$, and
\[
\operatorname{rank}\Lambda=2^{d-1}-d.
\]
\end{proof}

\begin{proof}[Proof of Corollary~\ref{cor:maximum}]
A partial Latin hyperrectangle in $\{0,1\}^d$ is an independent set in the
hypercube graph $Q_d$. A perfect matching gives $|L|\le2^{d-1}$, and the
bipartition classes $E_0,E_1$ attain equality.

Suppose $|L|=2^{d-1}$ and write $A=L\cap E_0$ and $B=L\cap E_1$. Then
$|E_1\setminus B|=|A|$. Since $L$ is independent, every edge from $A$ ends
in $E_1\setminus B$, so exactly $d|A|$ edges join these two sets. This is
the full degree sum $d|E_1\setminus B|$ of $E_1\setminus B$; hence there is
no edge between $E_0\setminus A$ and $E_1\setminus B$. There is also no edge
between $A$ and $B$. Thus $A\cup(E_1\setminus B)$ has no edge to its
complement. Since $Q_d$ is connected, this set is empty or all of $Q_d$.
Accordingly, $L=E_1$ or $L=E_0$, and the result follows from
Theorem~\ref{thm:main}.
\end{proof}

\section{The \texorpdfstring{$2\times2\times2$}{2 x 2 x 2} quartic}

Under the identification $x_k=i_k-1$, the even parity class is Ribot's
support
\[
L_2=\{(1,1,1),(1,2,2),(2,1,2),(2,2,1)\}.
\]
Take
\[
\Omega=\{\alpha,\beta,\gamma,\delta\},
\quad
\alpha=(1,1,1),
\quad
\beta=(1,1,-1),
\quad
\gamma=(1,-1,1),
\quad
\delta=(-1,1,1).
\]
The relation lattice is generated by
\[
\alpha-\beta-\gamma-\delta=0.
\]
With
\[
P=8p_\alpha,\qquad Q=8p_\beta,
\qquad R=8p_\gamma,\qquad S=8p_\delta,
\]
Theorem~\ref{thm:main} gives
\begin{equation}\label{eq:H}
X_{L_2}=\{T:H(T)=0\},
\qquad
H(T)=\Imag(P\overline Q\,\overline R\,\overline S).
\end{equation}

Ribot~\cite[Example~4.1]{Ribot} gives the Zariski closure associated with
$L_2$ by the quartic
\begin{equation}\label{eq:F}
F(T)=
\sum_{i,j,k\in\{1,2\}}(-1)^{i+j+k}t_{ijk}\,
g(t_{i+1,j,k},t_{i,j+1,k},t_{i,j,k+1},t_{i+1,j+1,k+1}),
\end{equation}
where shifted indices are read modulo $2$ and
\[
g(w,x,y,z)=2wxy+w^2z+x^2z+y^2z-z^3.
\]

\begin{proposition}\label{prop:quartic}
As polynomials in the eight real entries of $T$, one has $H=4F$.
Consequently,
\[
X_{L_2}=\{T:F(T)=0\}.
\]
\end{proposition}

\begin{proof}
The Fourier formula gives
\begin{align*}
P&=t_{111}-t_{122}-t_{212}-t_{221}
   +i(t_{112}+t_{121}+t_{211}-t_{222}),\\
Q&=t_{111}+t_{122}+t_{212}-t_{221}
   +i(-t_{112}+t_{121}+t_{211}+t_{222}),\\
R&=t_{111}+t_{122}-t_{212}+t_{221}
   +i(t_{112}-t_{121}+t_{211}+t_{222}),\\
S&=t_{111}-t_{122}+t_{212}+t_{221}
   +i(t_{112}+t_{121}-t_{211}+t_{222}).
\end{align*}
Substitution into \eqref{eq:H} and collection of degree-four monomials gives
$H=4F$; equivalently, the two sides have the same forty nonzero monomial
coefficients in $\Z[t_{ijk}:i,j,k\in\{1,2\}]$.
\end{proof}

Thus Conjecture~4.4 of~\cite{Ribot} holds for every maximum-cardinality
binary support. The argument does not treat non-parity binary supports or
formats with a mode of dimension at least three, and the conjecture remains
open in general.

\begin{thebibliography}{9}

\bibitem{BDHR}
A.~Boralevi, J.~Draisma, E.~Horobe\c{t}, and E.~Robeva,
\emph{Orthogonal and unitary tensor decomposition from an algebraic
perspective},
Israel J. Math. \textbf{222} (2017), no.~1, 223--260.

\bibitem{Ribot}
\'A.~Ribot,
\emph{Unifying singular value decompositions of tensors via aligned
orthogonality},
arXiv:2608.03202v1 [math.AG], 2026.

\end{thebibliography}

\end{document}